\begin{table}[!htbp]
\centering
\small
\setlength{\tabcolsep}{7pt}
\renewcommand{\arraystretch}{1.15}
\caption{Observed failure-free preissued-delivery latency (S6). Entries are median [minimum, maximum] of five run-level p95 latencies, in milliseconds.}
\label{tab:delivery-cost}
\begin{tabular}{@{}rrr@{}}
\toprule
Worker threads & ECRC & Ordinary\\
\midrule
1 & 33.6 [33.2, 33.9] & 33.1 [32.7, 33.6]\\
4 & 86.7 [58.6, 99.7] & 77.0 [66.8, 84.9]\\
8 & 204.0 [157.6, 231.0] & 166.0 [104.4, 203.5]\\
\bottomrule
\end{tabular}
\par\smallskip
\begin{minipage}{\linewidth}\small
Each thread/client condition has five runs of 64 timed actions at capacity 64: 30 runs and 1,920 actions overall, plus 30 separate untimed warm-ups. Latency covers arm to receipt, excluding permit issuance and model inference. Ranges describe the five observed runs; they are not confidence intervals. The synthetic loopback workload uses the stated single-host environment and is separate from the 90 recovery runs.
\end{minipage}
\end{table}
